\begin{abstract}

    \begin{center}
        \large{% Единое место для указания названия работы
Извлечение и анализ конечных автоматов из описаний аппаратуры на языке SystemVerilog
} \\
    \large\textit{Вязовцев Андрей Викторович} \\[1 cm]

Конечные автоматы (англ. Finite State Machines, FSM) являются широко используемым шаблоном проектирования
в языках описания аппаратуры (англ. Hardware Description Languages, HDL) для моделирования управляющей
логики. Несмотря на свою простоту, реализации FSM подвержены различным классам ошибок проектирования,
которые трудно обнаружить с помощью традиционных методов верификации на основе моделирования. В данной
работе представлен SVAN-FSM --- инструмент для автоматического извлечения и статического анализа конечных
автоматов из проектов на SystemVerilog. Инструмент работает на уровне списка логических элементов, строит
графовые представления и выявляет недостижимые состояния, <<мертвые>> переходы, блокировки (deadlocks)
и <<живые>> блокировки (livelocks). Также он обнаруживает взаимодействующие FSM для выявления межмодульных
зависимостей управления и поддерживает параметризованные проекты, анализируя каждую развернутую инстанцию
независимо. Результаты анализа сопоставляются с исходным кодом для повышения интерпретируемости.
Инструмент был протестирован на нескольких открытых проектах RISC-V и системах, критичных с точки
зрения безопасности, продемонстрировав свою эффективность на проектах промышленного масштаба.

    \vfill

    \textbf{Abstract} \\[1 cm]

Finite state machines (FSMs) are a widely used design pattern in hardware description languages (HDLs)
for modeling control logic. Despite their simplicity, FSM implementations are prone to various classes
of design errors that are difficult to detect using conventional simulation-based verification techniques.
This paper presents SVAN-FSM, a tool for automatic extraction and static analysis of FSMs from
SystemVerilog designs. The tool operates on RTL netlists, builds graph-based representations, and detects
unreachable states, dead transitions, deadlocks, and livelocks. It also identifies interacting FSMs
to expose cross-module control dependencies and handles parameterized designs by analyzing each elaborated
instance independently. Analysis results are mapped back to the source code to improve interpretability.
The tool has been evaluated on multiple open-source RISC-V designs and security-critical projects,
demonstrating its effectiveness on industrial-scale designs.

    \end{center}

\end{abstract}
